% =============================================================================
% CAPÍTULO 8 — GLOSARIO DE TÉRMINOS
% =============================================================================
\chapter{Glosario de Términos}
\label{ch:glosario}
\resumenCapitulo{Este capítulo define los términos técnicos y propios de la plataforma
empleados a lo largo del manual, para asegurar una comprensión uniforme por parte de
todas las personas usuarias.}

\begin{description}[style=unboxed,leftmargin=0pt,itemsep=4pt]
    \item[\textcolor{colorPrimario}{\textbf{Backoffice}}] Panel de administración de la
    plataforma, reservado al perfil Administrador, desde el cual se gestiona la
    operación global del sistema.

    \item[\textcolor{colorPrimario}{\textbf{CV Studio}}] Herramienta del perfil
    Profesional para crear, editar y exportar currículums a partir de plantillas en
    formato LaTeX o Markdown, con apoyo de un asistente de inteligencia artificial.

    \item[\textcolor{colorPrimario}{\textbf{Dashboard}}] Pantalla principal o panel de
    control de un perfil, que resume la información y los accesos más relevantes.

    \item[\textcolor{colorPrimario}{\textbf{Hard Skill (habilidad técnica)}}]
    Competencia técnica concreta y medible (por ejemplo, un lenguaje de programación o
    una tecnología), asociada a un nivel de dominio.

    \item[\textcolor{colorPrimario}{\textbf{Soft Skill (habilidad blanda)}}]
    Competencia interpersonal o de comportamiento (por ejemplo, trabajo en equipo o
    liderazgo), que se documenta con un caso de éxito.

    \item[\textcolor{colorPrimario}{\textbf{Match IA}}] Coincidencia automática entre
    una búsqueda en lenguaje natural y los perfiles disponibles, calculada por la
    inteligencia artificial de la plataforma.

    \item[\textcolor{colorPrimario}{\textbf{Pipeline}}] Flujo de selección del perfil
    Reclutador, organizado en columnas (Guardados, Revisión, Entrevista, Oferta y
    Contratado) por las que avanzan los candidatos.

    \item[\textcolor{colorPrimario}{\textbf{Portafolio}}] Conjunto de proyectos,
    experiencia laboral, formación académica y habilidades que conforman la presencia
    profesional de una persona en EthosHub.

    \item[\textcolor{colorPrimario}{\textbf{RBAC}}] \textit{Role-Based Access Control};
    modelo de control de acceso basado en roles que determina las funciones disponibles
    según el perfil de la cuenta.

    \item[\textcolor{colorPrimario}{\textbf{Responsive (diseño adaptable)}}] Capacidad
    de la interfaz para ajustarse automáticamente al tamaño de la pantalla, permitiendo
    su uso cómodo en teléfonos y tabletas.

    \item[\textcolor{colorPrimario}{\textbf{Top Skills}}] Sección destacada del
    portafolio donde se muestran las habilidades técnicas que la persona usuaria ha
    marcado como prioritarias.

    \item[\textcolor{colorPrimario}{\textbf{Ranking}}] Clasificación de talento que
    ordena a los profesionales según su competencia, con posibilidad de filtrar por
    área (Frontend, Backend, Data, Infraestructura, Mobile).

    \item[\textcolor{colorPrimario}{\textbf{Insight de la semana}}] Resumen automático
    del desempeño del reclutador generado por la plataforma en su dashboard.

    \item[\textcolor{colorPrimario}{\textbf{Modal}}] Ventana emergente que aparece sobre
    la pantalla actual para introducir datos o confirmar una acción sin abandonar el
    contexto.

    \item[\textcolor{colorPrimario}{\textbf{Toggle (interruptor)}}] Control de dos
    estados (activado/desactivado) empleado, por ejemplo, para definir la visibilidad
    de un proyecto.

    \item[\textcolor{colorPrimario}{\textbf{Administrador}}] Perfil con privilegios de
    gestión global sobre la plataforma: cuentas, roles, permisos, moderación de
    contenido y métricas del sistema.

    \item[\textcolor{colorPrimario}{\textbf{API}}] \textit{Application Programming
    Interface}; conjunto de reglas que permite que dos aplicaciones intercambien datos y
    se comuniquen entre sí.

    \item[\textcolor{colorPrimario}{\textbf{Autenticación}}] Proceso por el cual el
    sistema verifica la identidad de una persona usuaria mediante sus credenciales antes
    de concederle acceso.

    \item[\textcolor{colorPrimario}{\textbf{Avatar}}] Imagen o fotografía que representa
    a la persona usuaria en su perfil y en las interacciones dentro de la plataforma.

    \item[\textcolor{colorPrimario}{\textbf{Caché}}] Almacenamiento temporal de datos en
    el navegador o el servidor para acelerar la carga de páginas visitadas con
    frecuencia.

    \item[\textcolor{colorPrimario}{\textbf{Candidato}}] Profesional que ha sido
    incorporado por un reclutador a su \textit{pipeline} de selección para ser evaluado.

    \item[\textcolor{colorPrimario}{\textbf{Chat}}] Canal de mensajería interna que
    permite la comunicación directa y en tiempo real entre profesionales y reclutadores.

    \item[\textcolor{colorPrimario}{\textbf{Cookie}}] Pequeño archivo que el sitio
    almacena en el navegador para recordar la sesión y las preferencias de la persona
    usuaria.

    \item[\textcolor{colorPrimario}{\textbf{Credenciales}}] Conjunto formado por el
    correo electrónico y la contraseña que identifican a una cuenta al iniciar sesión.

    \item[\textcolor{colorPrimario}{\textbf{EthosHub}}] Plataforma web de gestión de
    talento y portafolios profesionales sobre la que trata este manual.

    \item[\textcolor{colorPrimario}{\textbf{Filtro}}] Criterio de búsqueda (área,
    habilidad, nivel, ubicación) que acota el listado de talento mostrado al reclutador.

    \item[\textcolor{colorPrimario}{\textbf{Frontend}}] Parte de una aplicación con la
    que interactúa directamente la persona usuaria; también designa un área del
    \textit{ranking} de talento.

    \item[\textcolor{colorPrimario}{\textbf{Backend}}] Parte de una aplicación que opera
    en el servidor y procesa la lógica y los datos; también un área del \textit{ranking}
    de talento.

    \item[\textcolor{colorPrimario}{\textbf{Inicio de sesión (login)}}] Procedimiento
    mediante el cual una persona usuaria accede a su cuenta introduciendo sus
    credenciales.

    \item[\textcolor{colorPrimario}{\textbf{JavaScript}}] Lenguaje de programación que
    ejecuta el navegador para habilitar las funciones interactivas de la plataforma.

    \item[\textcolor{colorPrimario}{\textbf{Markdown}}] Lenguaje de marcado ligero para
    dar formato a texto, admitido por el CV Studio para la edición de currículums.

    \item[\textcolor{colorPrimario}{\textbf{Metadatos}}] Información descriptiva asociada
    a un elemento (por ejemplo, fecha de creación o autor) que no forma parte de su
    contenido principal.

    \item[\textcolor{colorPrimario}{\textbf{Métrica}}] Valor cuantitativo que mide el
    desempeño o la actividad dentro de la plataforma (visitas, contactos, contrataciones).

    \item[\textcolor{colorPrimario}{\textbf{Moderación}}] Revisión del contenido
    publicado y de los reportes recibidos, realizada por el Administrador para garantizar
    el cumplimiento de las normas.

    \item[\textcolor{colorPrimario}{\textbf{Notificación}}] Aviso que la plataforma envía
    a la persona usuaria para informarle de un evento relevante (nuevo mensaje, cambio de
    estado, etc.).

    \item[\textcolor{colorPrimario}{\textbf{Onboarding}}] Proceso guiado de bienvenida
    que orienta a la persona usuaria en sus primeros pasos tras crear la cuenta.

    \item[\textcolor{colorPrimario}{\textbf{Perfil}}] Tipo de cuenta (Profesional,
    Reclutador o Administrador) que determina las funciones disponibles según el modelo
    RBAC.

    \item[\textcolor{colorPrimario}{\textbf{Permiso}}] Autorización concreta para
    ejecutar una acción dentro de la plataforma, otorgada en función del perfil.

    \item[\textcolor{colorPrimario}{\textbf{Plantilla}}] Diseño predefinido que sirve de
    base para generar un currículum en el CV Studio.

    \item[\textcolor{colorPrimario}{\textbf{Profesional}}] Perfil orientado a gestionar
    un portafolio propio, mostrar habilidades y generar currículums.

    \item[\textcolor{colorPrimario}{\textbf{Reclutador}}] Perfil orientado a buscar,
    filtrar y contactar talento, y a gestionar su \textit{pipeline} de candidatos.

    \item[\textcolor{colorPrimario}{\textbf{Reporte}}] Notificación que una persona
    usuaria envía para alertar sobre contenido inadecuado, sujeto a revisión por
    moderación.

    \item[\textcolor{colorPrimario}{\textbf{Token}}] Código temporal que la plataforma
    genera para validar acciones sensibles, como la verificación de correo o el
    restablecimiento de contraseña.

    \item[\textcolor{colorPrimario}{\textbf{Verificación}}] Confirmación de que una
    dirección de correo es válida y pertenece a la persona usuaria, mediante un código o
    enlace enviado a su bandeja.

    \item[\textcolor{colorPrimario}{\textbf{Asistente IA}}] Funcionalidad de
    inteligencia artificial que sugiere y redacta contenido para el portafolio y el CV
    Studio.

    \item[\textcolor{colorPrimario}{\textbf{Proyecto}}] Elemento del portafolio que
    documenta un trabajo realizado, con descripción, tecnologías empleadas y evidencias
    asociadas.

    \item[\textcolor{colorPrimario}{\textbf{Evidencia}}] Recurso (imagen, enlace o
    archivo) que respalda un proyecto y permite a un reclutador verificar su autenticidad.

    \item[\textcolor{colorPrimario}{\textbf{Experiencia laboral}}] Registro de la
    trayectoria profesional de la persona usuaria, con cargo, institución y periodo.

    \item[\textcolor{colorPrimario}{\textbf{Educación}}] Sección del portafolio que
    recoge la formación académica: institución, título, área de estudio y periodo.

    \item[\textcolor{colorPrimario}{\textbf{Skill}}] Habilidad técnica o blanda
    registrada en el portafolio y asociada a un nivel de dominio.

    \item[\textcolor{colorPrimario}{\textbf{Stack}}] Conjunto de tecnologías y
    herramientas que una persona usuaria domina o que se emplean en un proyecto.

    \item[\textcolor{colorPrimario}{\textbf{Sello de validación}}] Distintivo que
    certifica que una habilidad técnica ha superado el proceso de validación de la
    plataforma.

    \item[\textcolor{colorPrimario}{\textbf{Validación técnica}}] Procedimiento por el
    cual se solicita y aprueba la verificación formal de una habilidad técnica (hard
    skill).

    \item[\textcolor{colorPrimario}{\textbf{Normalización de skills}}] Proceso
    administrativo que unifica nombres equivalentes de habilidades en el repositorio
    global para evitar duplicados.

    \item[\textcolor{colorPrimario}{\textbf{Repositorio de habilidades}}] Catálogo
    global y centralizado de habilidades que el Administrador mantiene para toda la
    plataforma.

    \item[\textcolor{colorPrimario}{\textbf{Favorito}}] Perfil que un reclutador guarda
    para su consulta posterior dentro de sus listas de seguimiento.

    \item[\textcolor{colorPrimario}{\textbf{Lista de seguimiento}}] Colección de
    perfiles favoritos que el reclutador organiza para dar seguimiento al talento de
    interés.

    \item[\textcolor{colorPrimario}{\textbf{Conexión}}] Vínculo establecido entre dos
    personas usuarias dentro de la red profesional de EthosHub.

    \item[\textcolor{colorPrimario}{\textbf{Filtro avanzado}}] Combinación de criterios
    técnicos (habilidad, nivel, área, ubicación) que refina los resultados de búsqueda
    de talento.

    \item[\textcolor{colorPrimario}{\textbf{Motor de búsqueda}}] Componente que localiza
    perfiles a partir de palabras clave o lenguaje natural y los ordena por relevancia.

    \item[\textcolor{colorPrimario}{\textbf{Espacio de trabajo}}] Entorno principal de
    un perfil tras iniciar sesión, desde el cual se accede a todas sus funciones.

    \item[\textcolor{colorPrimario}{\textbf{Panel de reclutamiento}}] Vista principal
    del perfil Reclutador que reúne búsqueda de talento, pipeline y métricas.

    \item[\textcolor{colorPrimario}{\textbf{Panel de administración}}] Vista principal
    del perfil Administrador para la gestión global del sistema; véase \textit{Backoffice}.

    \item[\textcolor{colorPrimario}{\textbf{Perfil público}}] Versión visible del
    portafolio que otras personas usuarias pueden consultar según la configuración de
    privacidad.

    \item[\textcolor{colorPrimario}{\textbf{Privacidad del perfil}}] Conjunto de ajustes
    que determinan qué información del portafolio es visible para terceros.

    \item[\textcolor{colorPrimario}{\textbf{Identidad del perfil}}] Datos que
    identifican a la persona usuaria (nombre, avatar, nivel profesional) en su perfil
    público.

    \item[\textcolor{colorPrimario}{\textbf{Nivel profesional}}] Categoría que indica el
    grado de experiencia de un profesional (por ejemplo, Senior o Ingeniero Senior).

    \item[\textcolor{colorPrimario}{\textbf{Suspensión de cuenta}}] Medida
    administrativa que inhabilita temporalmente el acceso de una cuenta a la plataforma.

    \item[\textcolor{colorPrimario}{\textbf{Alta de cuenta}}] Activación o
    restablecimiento del acceso de una cuenta por parte del Administrador.

    \item[\textcolor{colorPrimario}{\textbf{Auditoría}}] Revisión del historial de
    acciones y cambios realizados en las cuentas para garantizar trazabilidad y
    cumplimiento.

    \item[\textcolor{colorPrimario}{\textbf{Rol}}] Conjunto de permisos asociado a un
    perfil; el Administrador puede modificarlo en casos especiales.

    \item[\textcolor{colorPrimario}{\textbf{Métricas globales}}] Indicadores agregados
    del estado y la actividad de toda la plataforma, accesibles para el Administrador.

    \item[\textcolor{colorPrimario}{\textbf{Monitoreo}}] Seguimiento continuo del estado
    del sistema y de sus métricas para detectar incidencias.

    \item[\textcolor{colorPrimario}{\textbf{Alerta}}] Aviso prioritario sobre contenido
    reportado o eventos que requieren la atención del Administrador.

    \item[\textcolor{colorPrimario}{\textbf{Exportación}}] Acción de descargar datos del
    portafolio o un currículum en un formato externo (PDF, TXT o Markdown).

    \item[\textcolor{colorPrimario}{\textbf{Exportar PDF}}] Opción del CV Studio que
    genera el currículum como documento PDF listo para compartir o imprimir.

    \item[\textcolor{colorPrimario}{\textbf{Zona de peligro}}] Sección de la
    configuración que agrupa acciones irreversibles, como la eliminación de la cuenta.

    \item[\textcolor{colorPrimario}{\textbf{LaTeX}}] Sistema de composición tipográfica
    empleado por el CV Studio para generar currículums con alta calidad de formato.

    \item[\textcolor{colorPrimario}{\textbf{Plantilla de CV}}] Diseño predefinido del CV
    Studio (Ejecutivo Clásico, Moderno Pro, Minimalista Senior, Compacto Una Página)
    sobre el que se construye el currículum.

    \item[\textcolor{colorPrimario}{\textbf{Tema (claro/oscuro)}}] Esquema de color de
    la interfaz que la persona usuaria puede alternar según su preferencia visual.

    \item[\textcolor{colorPrimario}{\textbf{Densidad de la interfaz}}] Ajuste de
    espaciado de la vista (por ejemplo, modo compacto) para mostrar más o menos
    información en pantalla.

    \item[\textcolor{colorPrimario}{\textbf{Recordarme}}] Opción del inicio de sesión
    que mantiene la sesión activa en el navegador hasta cerrarla manualmente.

    \item[\textcolor{colorPrimario}{\textbf{Cierre de sesión (logout)}}] Acción que
    finaliza la sesión activa y exige volver a autenticarse para acceder de nuevo.

    \item[\textcolor{colorPrimario}{\textbf{Recuperación de credenciales}}] Proceso para
    restablecer la contraseña de una cuenta mediante un código enviado al correo.

    \item[\textcolor{colorPrimario}{\textbf{Multimedia}}] Recursos visuales (imágenes,
    capturas o vídeos) que se adjuntan a un proyecto como evidencia.

    \item[\textcolor{colorPrimario}{\textbf{Captura de pantalla}}] Imagen del estado de
    la interfaz utilizada en este manual para ilustrar los procedimientos.

    \item[\textcolor{colorPrimario}{\textbf{Destinatario}}] Persona usuaria a la que se
    dirige un mensaje o una comunicación por correo dentro de la plataforma.

    \item[\textcolor{colorPrimario}{\textbf{Bandeja de entrada}}] Espacio donde se
    reciben los mensajes y las notificaciones dirigidos a la persona usuaria.

    \item[\textcolor{colorPrimario}{\textbf{Estado de solicitud}}] Etapa en la que se
    encuentra una petición (por ejemplo, de validación técnica): pendiente, en revisión
    o resuelta.

    \item[\textcolor{colorPrimario}{\textbf{FAQ}}] \textit{Frequently Asked Questions};
    sección de preguntas frecuentes que resuelve las dudas más habituales de las
    personas usuarias.

    \item[\textcolor{colorPrimario}{\textbf{Soporte técnico}}] Servicio de asistencia
    al que la persona usuaria recurre para resolver incidencias o reportar fallos
    críticos.

    \item[\textcolor{colorPrimario}{\textbf{Fallo crítico}}] Incidencia grave que impide
    el uso normal de la plataforma y requiere atención prioritaria del soporte.

    \item[\textcolor{colorPrimario}{\textbf{Área técnica}}] Especialidad del talento
    (Frontend, Backend, Data, Infraestructura, Mobile) usada para clasificar y filtrar
    perfiles.

    \item[\textcolor{colorPrimario}{\textbf{Etapa del pipeline}}] Cada una de las
    columnas por las que avanza un candidato (Guardados, Revisión, Entrevista, Oferta y
    Contratado).

    \item[\textcolor{colorPrimario}{\textbf{Guardados}}] Primera etapa del pipeline,
    donde el reclutador reúne los perfiles de interés antes de evaluarlos.

    \item[\textcolor{colorPrimario}{\textbf{Revisión}}] Etapa del pipeline en la que el
    reclutador analiza en detalle el portafolio y las evidencias del candidato.

    \item[\textcolor{colorPrimario}{\textbf{Entrevista}}] Etapa del pipeline en la que el
    candidato es contactado y evaluado mediante una conversación.

    \item[\textcolor{colorPrimario}{\textbf{Oferta}}] Etapa del pipeline en la que se
    formula una propuesta formal de contratación al candidato.

    \item[\textcolor{colorPrimario}{\textbf{Contratado}}] Etapa final del pipeline que
    marca al candidato como incorporado tras aceptar la oferta.

    \item[\textcolor{colorPrimario}{\textbf{Vacante}}] Posición o puesto que el
    reclutador busca cubrir mediante el proceso de selección.

    \item[\textcolor{colorPrimario}{\textbf{Tarjeta de candidato}}] Componente visual del
    pipeline que resume los datos clave de un candidato y se desplaza entre etapas.

    \item[\textcolor{colorPrimario}{\textbf{Descarte}}] Acción de retirar a un candidato
    del proceso de selección cuando no continúa en el pipeline.

    \item[\textcolor{colorPrimario}{\textbf{Caso de éxito}}] Relato breve que documenta
    una habilidad blanda mediante una situación real y su resultado.

    \item[\textcolor{colorPrimario}{\textbf{Certificado}}] Documento que acredita una
    formación o titulación y puede adjuntarse a la sección de educación del portafolio.

    \item[\textcolor{colorPrimario}{\textbf{Promedio (GPA)}}] Calificación media obtenida
    en una formación académica, registrada opcionalmente en la educación del portafolio.

    \item[\textcolor{colorPrimario}{\textbf{Área de estudio}}] Disciplina o campo
    académico asociado a un registro de formación (por ejemplo, Ingeniería de Software).

    \item[\textcolor{colorPrimario}{\textbf{Tecnología}}] Lenguaje, herramienta o
    framework que se asocia a un proyecto o a una habilidad técnica del portafolio.

    \item[\textcolor{colorPrimario}{\textbf{Sección del CV}}] Cada bloque de un currículum
    en el CV Studio (datos personales, experiencia, educación, habilidades, proyectos).

    \item[\textcolor{colorPrimario}{\textbf{Vista previa}}] Representación en pantalla del
    currículum o del perfil tal y como se verá una vez generado o publicado.

    \item[\textcolor{colorPrimario}{\textbf{Compilación}}] Proceso del CV Studio que
    transforma el contenido y la plantilla en el documento final (por ejemplo, en PDF).

    \item[\textcolor{colorPrimario}{\textbf{Aplicar al editor}}] Acción que inserta en el
    documento la sugerencia generada por el asistente de inteligencia artificial.

    \item[\textcolor{colorPrimario}{\textbf{Exportar TXT}}] Opción del CV Studio que
    genera el currículum como texto plano sin formato.

    \item[\textcolor{colorPrimario}{\textbf{Exportar Markdown}}] Opción del CV Studio que
    genera el currículum en formato Markdown para su edición posterior.

    \item[\textcolor{colorPrimario}{\textbf{Sesión}}] Periodo de actividad autenticada de
    una persona usuaria desde que inicia hasta que cierra sesión.

    \item[\textcolor{colorPrimario}{\textbf{Contraseña}}] Clave secreta que, junto al
    correo, autentica a la persona usuaria y debe cumplir los requisitos de seguridad.

    \item[\textcolor{colorPrimario}{\textbf{Datos personales}}] Información identificativa
    de la persona usuaria que la plataforma trata conforme a su política de privacidad.

    \item[\textcolor{colorPrimario}{\textbf{Política de privacidad}}] Documento que
    explica cómo la plataforma recopila, usa y protege los datos de las personas usuarias.

    \item[\textcolor{colorPrimario}{\textbf{Silenciar perfil}}] Acción de moderación que
    limita la visibilidad o la actividad de un perfil sin suspenderlo por completo.

    \item[\textcolor{colorPrimario}{\textbf{Resolución}}] Cierre de un reporte o
    incidencia tras la decisión tomada por el Administrador.

    \item[\textcolor{colorPrimario}{\textbf{Incidencia}}] Cualquier suceso anómalo
    registrado en el sistema que requiere seguimiento o intervención.

    \item[\textcolor{colorPrimario}{\textbf{Log (registro)}}] Anotación cronológica de
    las acciones y eventos del sistema, empleada en tareas de auditoría.

    \item[\textcolor{colorPrimario}{\textbf{Pestaña}}] Elemento de la interfaz que
    organiza el contenido en vistas separadas dentro de una misma pantalla.

    \item[\textcolor{colorPrimario}{\textbf{Barra lateral (sidebar)}}] Menú vertical de
    navegación situado al costado de la pantalla para acceder a las secciones del perfil.

    \item[\textcolor{colorPrimario}{\textbf{Insignia (badge)}}] Distintivo visual que
    señala un estado o logro, como un sello de validación o un nivel alcanzado.
\end{description}
